\documentclass[11pt,a4paper]{article}
\usepackage[margin=26mm,headheight=15pt]{geometry}
\usepackage{fontspec}
\setmainfont{texgyrepagella-regular.otf}[BoldFont=texgyrepagella-bold.otf,ItalicFont=texgyrepagella-italic.otf,BoldItalicFont=texgyrepagella-bolditalic.otf]
\usepackage{amsmath,amssymb,unicode-math}
\setmathfont{texgyrepagella-math.otf}
\usepackage{graphicx,booktabs,longtable,array,tabularx,xcolor,enumitem,fancyhdr}
\usepackage[unicode,colorlinks=true,linkcolor=teal!60!black,citecolor=teal!60!black,urlcolor=teal!60!black,bookmarksnumbered=true]{hyperref}
\usepackage{bookmark,xurl}
\graphicspath{{../figures/}}
\setlength{\emergencystretch}{2em}
\setlength{\parskip}{3pt}
\clubpenalty=10000\widowpenalty=10000\displaywidowpenalty=10000
\raggedbottom
\renewcommand{\topfraction}{.9}\renewcommand{\bottomfraction}{.85}
\renewcommand{\textfraction}{.08}\renewcommand{\floatpagefraction}{.85}
\setcounter{topnumber}{2}\setcounter{bottomnumber}{2}\setcounter{totalnumber}{3}
\setlist{itemsep=2pt,topsep=4pt}
\renewcommand{\arraystretch}{1.16}\setlength{\tabcolsep}{4pt}
\newcolumntype{Y}{>{\raggedright\arraybackslash}X}
\newcolumntype{L}[1]{>{\raggedright\arraybackslash}p{#1}}
\pagestyle{fancy}\fancyhf{}
\fancyhead[L]{\small Three-Way reciprocal hierarchy}
\fancyhead[R]{\small Paper II / v0.1}\fancyfoot[C]{\thepage}
\newtheorem{theorem}{Theorem}[section]
\newtheorem{proposition}[theorem]{Proposition}
\newtheorem{lemma}[theorem]{Lemma}
\newtheorem{corollary}[theorem]{Corollary}
\newtheorem{remark}[theorem]{Remark}
\newenvironment{proof}{\par\noindent\textit{Proof.}\ }{\hfill\(\square\)\par\medskip}
\numberwithin{equation}{section}
\newcommand{\PP}{\mathbb P}\newcommand{\CC}{\mathbb C}\newcommand{\RR}{\mathbb R}\newcommand{\ZZ}{\mathbb Z}
\newcommand{\Mon}{\operatorname{Mon}}\newcommand{\Deck}{\operatorname{Deck}}
\newcommand{\ord}{\operatorname{ord}}\newcommand{\Res}{\operatorname{Res}}\newcommand{\Disc}{\operatorname{Disc}}
\newcommand{\fig}[3]{\begin{figure}[htbp]\centering\includegraphics[width=\linewidth]{#1.pdf}\caption{#2}\label{#3}\end{figure}}
\hypersetup{pdftitle={The Three-Way Reciprocal Hierarchy: From Quadratic Involutions to Elliptic Correspondences and Symmetric Galois Closures},pdfauthor={Mathematical continuation prepared for Hilmir F. Halldorsson with Codex assistance},pdfsubject={Reciprocal rational maps, symmetric monodromy, elliptic correspondence, modular curves; novelty-neutral draft}}
\begin{document}
\hypersetup{pageanchor=false}
\begin{titlepage}\centering\vspace*{14mm}
{\Large The Three-Way Reciprocal Hierarchy\par}\vspace{7mm}
{\LARGE\bfseries From Quadratic Involutions to Elliptic Correspondences and Symmetric Galois Closures\par}
\vspace{12mm}
{\large Mathematical continuation prepared for Hilmir F. Halldórsson\par}
\vspace{3mm}{with Codex assistance\par}
\vspace{7mm}{Paper II / Version 0.1 / 5 October 2026\par}
\vspace{9mm}
\begin{minipage}{.93\textwidth}
\textbf{Abstract.} We extend the reconstructed Three-Way coefficient-reversal formula to
\(F_n(y)=P_n(y)/[y^nP_n(1/y)]\).
For every degree \(n\ge2\), we prove that a nonempty Zariski-open subset has symmetric monodromy \(S_n\) and Galois-closure genus
\(1+\frac{n!}{2}(n-3)\), over the unchanged target of \(F_n\).
The generic genera in degrees two, three, and four are therefore \(0,1,13\).
Degree two is exceptional because its unique other sheet is a Möbius deck partner. For the cubic, the two other sheets form an irreducible quadratic correspondence whose normalization is elliptic. Its \(S_3\) action consists of translations by 3-torsion and translated inversions, and its alternating quotient is an explicit 3-isogeny. Forgetting reciprocal marking gives \(X_0(3)\); retaining the distinguished nonzero 2-torsion gives \(X_0(6)\). We derive coordinates, modular invariants, and the cubic degeneration strata, then record a bounded quartic checkpoint. A final comparison preserves an exact finite \(S_3\)-incidence correspondence with six existing Tri-Octagon measurement patches. The inspected shell, scaffold, and response definitions supply no independently justified continuous modular identification. The comparison ends at the finite correspondence.
\end{minipage}
\vfill
\begin{minipage}{.93\textwidth}\small
\textbf{Status.} Novelty-neutral mathematical draft. Standard covering, elliptic, and modular theory is distinguished from explicit family calculations and the project-specific comparison. No physics interpretation, submission, or publication is implied. Attribution and publication metadata remain subject to author review. Paper I and the source research reports are preserved unchanged.
\end{minipage}
\end{titlepage}
\hypersetup{pageanchor=true}\pagenumbering{roman}
{\small\setlength{\parskip}{0pt}\tableofcontents}\clearpage\pagenumbering{arabic}
\section{The hierarchy and the cover being studied}\label{sec:scope}
Paper I reconstructs the historical quadratic Three-Way formula, its reciprocal quotient, and its additional squaring cover \cite{paper1}. The present paper asks a different mathematical question: what survives when coefficient reversal is extended to arbitrary degree? We begin with
\begin{equation}\label{eq:family}
 P_n(y)=a_0y^n-\sum_{j=1}^{n-1}a_jy^{n-j}+1,
 \qquad Q_n(y)=y^nP_n(1/y),\qquad F_n=P_n/Q_n.
\end{equation}
The coefficients range over \(\CC\); the minus signs impose no positivity condition. All covers are reduced rational maps between smooth projective curves, in characteristic zero. ``Generic'' means a nonempty Zariski-open subset of the indicated coefficient space. A real coefficient or real plotting statement will be identified explicitly.

Throughout the main theorem, the base field is \(\CC(v)\), with \(v=F_n(y)\). The Galois closure is the smooth connected curve with the splitting field of the generic fiber polynomial as its function field. Its degree is the order of the monodromy group. Neither target inversion nor the substitution \(y=x^2\) is incorporated into that extension unless explicitly stated.

This convention resolves a potentially misleading comparison. The degree-two cover \(F_2\) is already Galois and has sphere source, so its own closure has genus zero. The Edwards genus-one curve of Paper I belongs to the larger degree-eight quotient/lift tower studied there. Changing the base quotient or adding a source cover changes the function-field extension. The two genus statements are compatible, and no identification of their elliptic curves is asserted.

The results here assemble two external exact-analysis reports \cite{hierarchy,cubic}. The basic species is standard: after Cayley coordinates these are odd rational functions, an order-two instance of the normal form for maps commuting with a rotation \cite[§2.1, Lemma 2]{msw}. A commuting automorphism in that dynamical usage need not fix the value of the map. We reserve \emph{deck transformation} for a source automorphism that fixes the target pointwise. General low-genus Galois closures of rational maps are also established territory \cite[§2 and Theorems 1.1--1.2]{pakovich}. Our purpose is an explicit, verified account of this coefficient chart, without a priority claim.

\subsection{Reduction precedes locks and special values}
Write \(P=P_n\), \(Q=Q_n\), \(G=\gcd(P,Q)\) monic, and
\begin{equation}\label{eq:reduced}
 N=P/G,\qquad D=Q/G,\qquad F=N/D.
\end{equation}
Since \(P(0)=1\) and \(Q\) is monic of degree \(n\), no common zero occurs at zero or infinity. For a nonconstant map,
\begin{equation}\label{eq:degree}
 \deg F=n-\deg G.
\end{equation}
In particular, \(a_0=0\) need not lower the map degree. The displayed unreduced fraction still has holes at common zeros, even though its reduced rational continuation is defined projectively there.

\begin{proposition}[Reciprocal divisors]\label{prop:reciprocal}
For every coefficient choice, as a rational identity,
\begin{equation}\label{eq:reciprocity}
 F(1/y)=1/F(y).
\end{equation}
For a nonconstant map and \(R(y)=1/y\),
\begin{equation}\label{eq:divisor}
 \ord_{R(a)}F=-\ord_aF.
\end{equation}
Thus surviving zeros and poles pair reciprocally with equal multiplicity, including zero and infinity. Common factors occur in reciprocal pairs, except at the fixed points \(\pm1\).
\end{proposition}
\begin{proof}
Reversal twice gives \(Q^*=P\). Hence \(P(1/y)/Q(1/y)=Q(y)/P(y)\); this remains true after reduction. Taking local orders proves the divisor statement. A common root and its reciprocal have common multiplicities interchanged by reversal.
\end{proof}

For \(a_0\ne0\), write \(P=a_0\prod_i(y-r_i)\). Then \(Q=a_0\prod_i(1-r_i y)\), and
\begin{align}\label{eq:resultant_general}
 \Res(P,Q)&=a_0^{2n}\prod_{i,j}(1-r_i r_j)\notag\\
 &=(-1)^na_0^{2n-2}P(1)P(-1)
       \prod_{i<j}(1-r_i r_j)^2.
\end{align}
Cancellation is therefore equivalent to a reciprocal root pair, allowing a repeated index at \(r_i=\pm1\). At \(a_0=0\), use the polynomial resultant or the gcd instead of this degree-\(n\) root parameterization. Repeated surviving zeros are ramified points over zero; their reciprocal poles have the same multiplicities. A double zero alone is simple ramification and does not force a drop in generic monodromy.

\subsection{Parity and the fixed-value fibers}
The correct homogeneous fiber equations are
\begin{equation}\label{eq:fibers}
 F=1:\ N-D=0,\qquad F=-1:\ N+D=0.
\end{equation}
Equivalently divide \(P-Q\) and \(P+Q\) by \(G\) first. Projective interpretation includes infinity; at a finite solution the reduced denominator must be nonzero. The equation \(F=1/F\) is exactly the union of these two fibers, since neither zero nor infinity is fixed by target inversion.

The form \(P-Q\) is anti-reciprocal and \(P+Q\) is reciprocal. Their forced factors are
\begin{center}
\begin{tabular}{lll}\toprule
Parity & \(P-Q\) & \(P+Q\)\\\midrule
Even \(n\) & \((y-1)(y+1)\) & no forced linear factor\\
Odd \(n\) & \(y-1\) & \(y+1\)\\\bottomrule
\end{tabular}
\end{center}
Consequently an uncanceled fraction has
\begin{equation}\label{eq:locks}
 F(1)=1,\qquad F(-1)=(-1)^n.
\end{equation}
There is a precise correction at a canceled fixed point.

\begin{proposition}[Lock continuation]\label{prop:lock}
If \(P\) vanishes to order \(m\) at \(\epsilon\in\{1,-1\}\), its reduced continuation satisfies
\begin{equation}\label{eq:lock_continuation}
 F(\epsilon)=\epsilon^n(-1)^m.
\end{equation}
\end{proposition}
\begin{proof}
Put \(y=\epsilon+h\). Since \(1/y-\epsilon=-h+O(h^2)\), reversal multiplies the leading local coefficient by \(\epsilon^n(-1)^m\). Dividing gives the stated value.
\end{proof}
Thus odd-order cancellation flips a lock. For example, at \((A,B)=(2,3)\), \(F_2\) reduces to \((2y-1)/(y-2)\), whose value at \(y=1\) is \(-1\). The original fraction has a hole there.

For \(a_0\ne0\), \(F(0)=1/a_0\) and \(F(\infty)=a_0\). If \(a_0=0\) and \(d_P=\deg P<n\), zero is a pole and infinity a zero, both of order \(n-d_P\). Finally the only constant maps are \(+1\) and \(-1\). They occur respectively when
\begin{align}\label{eq:constant}
 a_0=1,&\quad a_j=a_{n-j},\notag\\
 a_0=-1,&\quad a_j=-a_{n-j}.
\end{align}
For even \(n\), the second condition forces the middle coefficient to vanish. These statements follow by reversing \(P=cQ\), which forces \(c^2=1\).

\subsection{Critical pairing and Cayley oddness}
The reduced critical divisor is obtained by homogenizing
\begin{equation}\label{eq:wronskian}
 W_{\rm red}=N'D-ND',\qquad P'Q-PQ'=G^2W_{\rm red}
\end{equation}
to degree \(2\deg F-2\). This includes ramified poles and infinity. A zero of order \(m\) in the divisor has local degree \(m+1\). Direct differentiation of reversal gives
\begin{equation}\label{eq:W_reversal}
 W(1/y)=y^{2-2n}W(y),\qquad W=P'Q-PQ'.
\end{equation}
Hence critical points and values pair as \((c,F(c))\leftrightarrow(1/c,1/F(c))\).

Let
\begin{equation}\label{eq:cayley}
 \mathcal C(y)=\frac{y-1}{y+1},\qquad
 z=\mathcal C(y),\qquad H=\mathcal C\circ F\circ\mathcal C^{-1}.
\end{equation}
Then \(\mathcal C(1/y)=-\mathcal C(y)\), so
\begin{equation}\label{eq:odd}
 H(-z)=-H(z),\qquad H(z)=z\Psi(z^2),\quad \Psi\in\CC(T).
\end{equation}
Here \(T\) is an independent variable: \(H(z)/z\) belongs to the fixed field \(\CC(z^2)\), so \(\Psi\) is a rational function evaluated at \(z^2\). The generic forms in the uncanceled coefficient chart are
\begin{equation}\label{eq:odd_types}
 \begin{array}{ll}
 n=2m+1:& H(z)=zU_m(z^2)/V_m(z^2),\\
 n=2m:& H(z)=zU_{m-1}(z^2)/V_m(z^2).
 \end{array}
\end{equation}
The leading and constant coefficients displayed here are generically nonzero. The odd-degree component fixes zero and infinity; the even-degree component sends both to zero. Cancellation may change component and degree.

At a source fixed point of reciprocity the target must also be fixed by inversion. In local coordinates that turn both involutions into sign change, the map is odd. Its local degree is therefore odd. In particular \(y=\pm1\) cannot be simple critical points, although local degree three or higher is possible exceptionally.

For real coefficients, conjugation and reciprocity imply \(|F(y)|=1\) on \(|y|=1\), wherever the raw fraction is defined. A surviving zero or pole cannot lie on that circle: conjugation preserves its divisor sign while inversion reverses it. The reduced nonconstant map therefore extends across raw holes on the circle with unit-modulus values. This unit-circle statement concerns a complex source contour, not the real plotting interval \(y\ge0\).

\section{Generic symmetric monodromy and closure genus}\label{sec:generic}
\begin{theorem}\label{thm:generic}
For every \(n\ge2\), the family \eqref{eq:family} contains a nonempty Zariski-open subset of degree-\(n\) maps with \(2n-2\) simple critical points and distinct critical values. On this subset,
\begin{equation}\label{eq:main}
 \boxed{\Mon(F_n)=S_n,\qquad
 g(\widetilde F_n)=1+\frac{n!}{2}(n-3).}
\end{equation}
For \(n\ge3\), the deck group is trivial and the residual equal-value correspondence has degree \(n-1\) over \(\CC(y)\) and is irreducible. These statements concern the unchanged target \(v=F_n(y)\).
\end{theorem}

\subsection{A construction proving the generic locus is nonempty}
It is not enough to invoke generic rational maps: coefficient reversal imposes a constraint, and one must show that this constraint does not force repeated branch values. The following construction does so inside the odd coordinate chart.

Start with \(H_1(z)=z\) for odd degrees and \(H_2(z)=z/(z^2+1)\) for even degrees. The first has no critical points; the second has simple critical points \(\pm1\), with values \(\pm1/2\). Both are unramified at zero and infinity. Their finite poles, if any, are simple, and their asymptotic behavior is respectively \(cz\) and \(c/z\), with \(c\ne0\).

Suppose \(H_d\) has these properties, with all critical points simple and all critical values distinct. Choose a finite \(b\ne0\) such that \(\pm b\) are neither poles nor critical points and \(H_d(b)\) is finite, nonzero, and different from every old critical value and its negative. Only finitely many choices are forbidden. For small \(\varepsilon\ne0\), define
\begin{equation}\label{eq:insertion}
 H_{d+2}(z)=H_d(z)+\varepsilon\frac{z}{z^2-b^2}.
\end{equation}
It remains odd, adds simple poles at \(\pm b\), and has degree \(d+2\). The old poles keep their nonzero residues. The numerator cannot cancel a new pole because its residue is \(\varepsilon/2\). The behavior at infinity retains its type and its nonzero leading coefficient for sufficiently small \(\varepsilon\).

Write \(a=H_d'(b)\ne0\) and \(z=b+\xi\). The new pole contributes \(\varepsilon/(2\xi)+O(\varepsilon)\). Solving the derivative equation by a local square-root expansion gives two new critical points and their values:
\begin{align}\label{eq:local_insertion}
 z&=b\pm\sqrt{\varepsilon/(2a)}+O(\varepsilon),\notag\\
 H_{d+2}(z)&=H_d(b)\pm\sqrt{2\varepsilon a}+O(\varepsilon).
\end{align}
For a rigorous local parameter, put \(\varepsilon=t_0^2\) and \(\xi=t_0\zeta\); after multiplication by \(\zeta^2\), the limiting derivative equation has two simple roots \(\zeta=\pm(2a)^{-1/2}\). The implicit-function theorem gives the displayed branches. Their critical points are simple. Oddness supplies the two opposite branches near \(-b\). Since \(H_d(b)\ne0\), these four values are distinct for sufficiently small nonzero \(t_0\) and avoid the old values.

Old simple critical points persist by the same theorem. Nonzero local linear coefficients at zero and infinity persist. We have exhibited
\((2d-2)+4=2(d+2)-2\) simple critical points. Riemann--Hurwitz leaves no additional ramification. Induction therefore constructs the required odd maps in every degree.

To return to the precise coefficient chart, conjugate by \(\mathcal C^{-1}\). A generic nonzero target scaling of \(H\) avoids the condition that the resulting numerator vanish at \(y=0\); it preserves oddness and branch simplicity. Reciprocity makes the resulting homogeneous numerator and reversed denominator proportional. At \(y=1\), the value is \(+1\), with no cancellation, so the proportionality sign is \(+1\). Scaling the pair to make the numerator constant term one gives \(Q=P^*\) exactly. The parity types in \eqref{eq:odd_types} give the required degree and value at \(-1\). Thus these are witnesses in \eqref{eq:family}, not in an unrelated larger family.

Coprimeness, simple critical divisor, and distinct critical values are open algebraic conditions: their failures are detected by homogeneous resultants and discriminants. The coefficient space \(\CC^n\) is irreducible. Existence of a witness in each degree consequently proves that the desired locus is nonempty and Zariski open.

\subsection{Monodromy, partners, and genus}
At a simple branch value precisely one sheet pair is exchanged. All branch cycles are therefore transpositions, and they generate the monodromy group. Connectedness of the source makes that action transitive. A transitive group generated by transpositions is \(S_n\): form the graph whose edges are those transpositions; it is connected, and transpositions along a connected graph generate every permutation.

Deck permutations commute with monodromy on a regular fiber. The centralizer of the natural \(S_n\) action is trivial for \(n\ge3\). After choosing one sheet, its stabilizer is \(S_{n-1}\), which acts transitively on the other \(n-1\) sheets. The residual equal-value correspondence is therefore irreducible over the chosen source field. A rational partner \(M\) would satisfy \(F_n\circ M=F_n\); multiplicativity of degree forces \(\deg M=1\), making it a Möbius deck transformation. None exists generically for \(n\ge3\).

Finally the Galois closure has degree \(n!\). It has \(2n-2\) branch values, each with inertia order two. Galois Riemann--Hurwitz yields
\begin{equation}\label{eq:RH_general}
 2g-2=n!\left[-2+(2n-2)\left(1-\frac12\right)\right]
       =n!(n-3).
\end{equation}
This proves Theorem~\ref{thm:generic}. The construction proves all degrees; no extrapolation from the two numerical witnesses is used.

\subsection{Why the quadratic member is exceptional}\label{sec:quadratic}
A separable degree-two extension is automatically Galois. Geometrically there is exactly one other point in a generic fiber, so the partner cannot be permuted with competing choices. It is a rational automorphism of the sphere and hence Möbius. For the uncanceled quadratic Three-Way map,
\begin{equation}\label{eq:quadratic_partner}
 T(y)=\frac{(A+1)y-B}{By-(A+1)},\qquad F_2(T(y))=F_2(y).
\end{equation}
The degree-two condition is \(A\ne1\) and \((A+1)^2-B^2\ne0\). The latter is also the nonzero determinant condition for \(T\). Reciprocity \(R(y)=1/y\) normalizes the unique nonidentity deck element, so it commutes with \(T\). They are distinct and generate \(V_4\). Its subgroup \(\{1,T\}\) fixes \(F_2\); its other two elements invert \(F_2\). Thus \(V_4\) fixes the further target quotient, for example \(F_2+1/F_2\), while \(\Deck(F_2)=C_2\).

The following table always concerns \(F_n\) itself, after cancellation and on its generic locus. The lock column assumes no cancellation at the indicated points.
\begin{table}[htbp]\centering\small
\begin{tabularx}{\linewidth}{L{.25\linewidth}YYY}\toprule
Property & \(n=2\) & \(n=3\) & \(n=4\)\\\midrule
Degree & 2 & 3 & 4\\
Reciprocity & \(F(1/y)=1/F(y)\) & same & same\\
Other-sheet equation & linear, rational partner & irreducible quadratic & irreducible cubic\\
Deck group & \(C_2\) & trivial & trivial\\
Monodromy & \(C_2=S_2\) & \(S_3\) & \(S_4\)\\
Closure degree; genus & 2; 0 & 6; 1 & 24; 13\\
Forced \(+1\) locks & \(y=\pm1\) & \(y=1\) & \(y=\pm1\)\\
Selected exceptions & cancellation / constant & cyclic, nodal, cancellation & decomposable, cyclic, \(V_4\), cancellation\\\bottomrule
\end{tabularx}
\caption{Generic hierarchy over the unchanged rational-map target. Even parity restores a lock at \(-1\), but does not restore the quadratic partner mechanism.}\label{tab:hierarchy}
\end{table}
\fig{01_hierarchy}{Degrees and closure genera refer to three distinct extensions over their own unchanged targets. Paper I's elliptic curve is associated with the separate larger quadratic tower. The lower row records the residual other-sheet degree.}{fig:hierarchy}

\section{The cubic equal-value correspondence}\label{sec:cubic}
Let
\begin{equation}\label{eq:cubic_original}
 F_3(y)=\frac{Ay^3-By^2-Cy+1}{y^3-Cy^2-By+A}=\frac{P(y)}{Q(y)}.
\end{equation}
The elementary fiber and cancellation identities are
\begin{align}\label{eq:cubic_fibers}
 P-Q&=(y-1)[(A-1)(y^2+y+1)+(C-B)y],\notag\\
 P+Q&=(y+1)[(A+1)(y^2-y+1)-(B+C)y],\\
 \Res(P,Q)&=(A-B-C+1)(A+B-C-1)H_3^2,\label{eq:cubic_resultant}\\
 H_3&=A^2+AC-B-1.\notag
\end{align}
The two linear factors of the resultant correspond to cancellation at \(\pm1\); the square corresponds to reciprocal-pair cancellation, including intersections with the fixed-point strata. Zeros are the roots of \(P/G\), with reciprocal poles. When \(A\ne0\),
\begin{equation}\label{eq:cubic_zero_disc}
 \Disc(P)=B^2C^2+4AC^3+4B^3+18ABC-27A^2.
\end{equation}
For a real uncanceled cubic, a positive value gives three distinct real zeros and a negative value one real zero and a conjugate nonreal pair. A zero discriminant signals repetition. These familiar real-cubic facts do not decide which zeros lie in \(y\ge0\).

\subsection{Exact factorization in the original source}
Put \(a=B-AC\), \(b=C-AB\), \(c=A^2-1\). Direct expansion gives
\begin{equation}\label{eq:original_factor}
 F_3(u)-F_3(y)=\frac{(u-y)\mathcal Q_y(u)}{Q(u)Q(y)},
\end{equation}
where
\begin{align}\label{eq:original_residual}
 \mathcal Q_y(u)={}&(ay^2+by+c)u^2\notag\\
 &+[by^2+(c+B^2-C^2)y+b]u+cy^2+by+a.
\end{align}
The quadratic describes the two other sheets. Its leading coefficient can vanish at a particular source point; this is a projective phenomenon, not a global rational splitting. On the degree-three locus, any rational root over \(\CC(y)\) would give a Möbius deck transformation. Since the order of a deck group divides the cover degree, one nonidentity deck element forces a cyclic Galois cover of degree three. The exceptional splitting criterion will therefore coincide with the cyclic criterion below.

\subsection{Odd coordinates and the discriminant curve}
In \(z=(y-1)/(y+1)\) and \(w=(F_3-1)/(F_3+1)\), set
\begin{equation}\label{eq:pqrs}
 \begin{aligned}
 p&=A+B-C-1,&q&=3(A-1)-B+C,\\
 r&=3(A+1)+B+C,&s&=A+1-B-C.
 \end{aligned}
\end{equation}
Then
\begin{equation}\label{eq:odd_cubic}
 w=H(z)=\frac{z(pz^2+q)}{rz^2+s},\qquad r+s-p-q=8.
\end{equation}
The last relation fixes the common scalar in these coefficients. The homogeneous numerator and denominator are \(pZ^3+qZT^2\) and \(rZ^2T+sT^3\). Consequently
\begin{equation}\label{eq:degree3}
 \deg H=3\quad\Longleftrightarrow\quad ps(ps-qr)\ne0,
 \qquad ps-qr=-8H_3.
\end{equation}
After removing \(v-z\) from the equality \(H(v)=H(z)\), the residual relation is
\begin{equation}\label{eq:K}
 \mathcal K_z(v)=p(rz^2+s)v^2+(ps-qr)zv+s(pz^2+q)=0.
\end{equation}
Its discriminant is
\begin{equation}\label{eq:D}
 \mathcal D(z)=-4p^2rsz^4+(q^2r^2-6pqrs-3p^2s^2)z^2-4pqs^2.
\end{equation}
Adjoining its square root recovers the second sheet via
\begin{equation}\label{eq:eta_incidence}
 \eta=2p(rz^2+s)v+(ps-qr)z.
\end{equation}
The third sheet is then rational in the first two, since the fiber cubic is
\begin{equation}\label{eq:fiber_cubic}
 pT^3-wrT^2+qT-ws=0,
 \qquad u=\frac rp w-z-v.
\end{equation}
Thus the smooth normalization of
\begin{equation}\label{eq:E}
 E:\quad\eta^2=\mathcal D(z)
\end{equation}
has the full splitting field over \(\CC(w)\) as its function field, whenever \(\mathcal K\) is irreducible. It is not merely a curve introduced by analogy.

\subsection{Critical points and branch values}
Differentiation and a cubic discriminant give respectively
\begin{align}\label{eq:J}
 J_z&=prz^4+(3ps-qr)z^2+qs,\\
 \mathcal B(w)&=-4r^3sw^4+(q^2r^2+18pqrs-27p^2s^2)w^2-4pq^3.
 \label{eq:branch_quartic}
\end{align}
Interpret both as binary quartics, retaining roots at infinity. The first is the critical divisor; the second is \(\Disc_T(pT^3-wrT^2+qT-ws)\), whose zeros are the branch values. In original target coordinates, the latter are \((1+w)/(1-w)\). Exact discriminants are
\begin{align}\label{eq:cubic_discriminants}
 \Disc(J_z)&=16pqrs(ps-qr)^2(9ps-qr)^2,\notag\\
 \Disc(\mathcal D)&=256p^3qrs^3(ps-qr)^6(9ps-qr)^2.
\end{align}
In particular the smooth elliptic locus is
\begin{equation}\label{eq:elliptic_locus}
 pqrs(ps-qr)(9ps-qr)\ne0.
\end{equation}
There the quartic \(\mathcal D\) has four distinct roots. A double cover of the sphere branched at four points has genus one, so \(E\) is a genus-one curve. The map \(E\to\PP^1_w\) has degree six and group \(S_3\), consistently with Theorem~\ref{thm:generic}.

\subsection{Six root orderings and the elliptic action}\label{sec:action}
Over a regular \(w\), a point of \(E\) specifies an ordered pair \((z,v)\) of distinct roots, or equivalently the full ordering \((z,v,u)\). There are six such points. Compactification also includes the limiting configurations with coincident roots; deleting them would lose ramification.

The field automorphisms
\begin{equation}\label{eq:generators}
 \tau(z,v,u)=(v,u,z),\qquad \iota(z,v,u)=(z,u,v)
\end{equation}
preserve \(w\) and satisfy
\(\tau^3=\iota^2=1\), \(\iota\tau\iota=\tau^{-1}\).
In quartic coordinates, \(\iota(z,\eta)=(z,-\eta)\). An explicit formula for \(\tau\) follows by solving \eqref{eq:eta_incidence} for \(v\), using \eqref{eq:fiber_cubic} for \(u\), and taking
\((v,,2p(rv^2+s)u+(ps-qr)v)\) as the new quartic coordinates.

Choose one branch point \(a\) of \(\mathcal D\), and let \(O=(a,0)\) be the elliptic origin. This choice is over \(\CC\); a real origin is not assumed. Then \(\iota=[-1]\).

\begin{proposition}\label{prop:torsion_action}
There is a nonzero point \(P\in E[3]\) such that \(\tau=T_P\). The six deck automorphisms of \(E\to\PP^1_w\) are
\begin{equation}\label{eq:six_autos}
 1,\ T_P,\ T_{-P},\ [-1],\ T_P[-1],\ T_{-P}[-1].
\end{equation}
\end{proposition}
\begin{proof}
Every inertia group has order two, so an element of order three has no fixed point. Alternatively, a fixed ordering under \(\tau\) would require a triple root, absent on \eqref{eq:elliptic_locus}. Every elliptic automorphism is \(T_R a\), with \(a\) fixing the origin. If \(a\ne1\), the nonzero isogeny \(1-a\) is surjective, giving a solution of \((1-a)X=R\) and thus a fixed point. Hence \(\tau\) is a translation; its exact order implies \(P\ne O\) and \(3P=O\). The relations in \eqref{eq:generators} give the remaining elements.
\end{proof}
The fixed points of the three transpositions solve \(2X=kP\), for \(k=0,1,-1\). Each transposition has four fixed points. Only the selected \([-1]\) is origin-fixing in this description. An abstract group count alone would not establish these translations or fixed-point sets.

\fig{02_incidence}{A regular cubic fiber has three sheets and six ordered pairs. Root-label permutations act on this six-element set; choosing the first entry gives the three-sheet quotient. Position permutations give the commuting regular deck action. The elliptic normalization includes the branch limits as well as these regular fibers.}{fig:incidence}

\section{Elliptic equations, isogeny, and marked moduli}\label{sec:modular}
\subsection{An explicit Weierstrass model}
For compact notation put
\begin{equation}\label{eq:UV}
 U=ps,\quad V=qr,\quad
 \alpha=-4p^2rs,\quad\beta=V^2-6UV-3U^2,\quad\gamma=-4pqs^2.
\end{equation}
Then \(\mathcal D(z)=f(z)=\alpha z^4+\beta z^2+\gamma\). Let
\begin{equation}\label{eq:IJ}
 I=\beta^2+12\alpha\gamma,\qquad J=72\alpha\beta\gamma-2\beta^3.
\end{equation}
With \(a\) and \(O\) chosen as above, the birational transformation
\begin{equation}\label{eq:weierstrass_map}
 X=\frac{f'(a)}{4(z-a)}+\frac{f''(a)}{24},\qquad
 Y=-\frac{f'(a)\eta}{8(z-a)^2}
\end{equation}
gives
\begin{equation}\label{eq:weierstrass}
 Y^2=X^3-\frac I{48}X-\frac J{1728}.
\end{equation}
The inverse is \(z=a+f'(a)/[4(X-f''(a)/24)]\), followed by
\(\eta=-8Y(z-a)^2/f'(a)\). Substitution using \(f(a)=0\) proves the identity; the exceptional affine denominators are filled on the smooth projective curves. The invariants in this convention are
\begin{align}\label{eq:invariants}
 c_4&=I,\qquad c_6=J/2,\notag\\
 \Delta_E&=\frac{\alpha\gamma(\beta^2-4\alpha\gamma)^2}{16}
   =p^3qrs^3(ps-qr)^6(9ps-qr)^2,\notag\\
 j(E)&=\frac{16(\beta^2+12\alpha\gamma)^3}
                 {\alpha\gamma(\beta^2-4\alpha\gamma)^2}.
\end{align}
In particular \(\Disc(f)=256\Delta_E\), which reconciles the quartic and Weierstrass discriminant conventions.

For a consistent orientation of \(\tau\), its translation point in this model is
\begin{equation}\label{eq:P}
 P=\left(\frac{(V+3U)^2}{12},-\frac{U(V-U)^2}{2}\right).
\end{equation}
It lies on \eqref{eq:weierstrass}, and the tangent doubling formula gives \(2P=-P\). Its ordinate is nonzero on \eqref{eq:elliptic_locus}. Reversing the cyclic generator changes its sign. The link to the actual incidence action, rather than just an independently found torsion point, is made explicit below.

\subsection{Reciprocity selects nonzero 2-torsion}
Simultaneous reciprocity on the three roots lifts as
\begin{equation}\label{eq:rho}
 \rho(z,v,u)=(-z,-v,-u),\qquad \rho(z,\eta)=(-z,-\eta).
\end{equation}
It commutes with every root-position permutation and sends \(w\) to \(-w\). Its only possible fixed source coordinates are zero and infinity. It swaps the two points above zero, since \(\gamma\ne0\); at infinity it exchanges the points distinguished by \(\eta/z^2=\pm\sqrt\alpha\). It is therefore fixed-point-free. The argument of Proposition~\ref{prop:torsion_action} gives
\begin{equation}\label{eq:Q}
 \begin{gathered}
 \rho=T_Q,\qquad Q\in E[2],\quad Q\ne O,\\
 Q=(-a,0)\ \text{on the quartic},\qquad
 Q=(-\beta/6,0)\ \text{on \eqref{eq:weierstrass}}.
 \end{gathered}
\end{equation}
The other lift \((z,\eta)\mapsto(-z,\eta)\) is \(T_Q[-1]\), has fixed points, and does not commute with the entire \(S_3\) action. Thus the commuting lift characterizes the selected 2-torsion. The translations \(\langle P,Q\rangle\) form a cyclic group of order six, and together with inversion generate \(S_3\times C_2\). Only its \(S_3\) subgroup fixes \(w\).

\subsection{The alternating quotient is a 3-isogeny}\label{sec:isogeny}
Quotienting by \(\langle\iota\rangle\) recovers the original \(z\)-sphere. Quotienting instead by the normal translation group \(\langle\tau\rangle=C_3\) gives a genus-one curve \(E'\), since that action is free. With origin \(O'=\phi(O)\), its quotient map is an isogeny of degree three.

The model and map are explicit:
\begin{equation}\label{eq:isogeny}
 \boxed{E':\chi^2=\mathcal B(w),\qquad
 \phi(z,\eta)=\left(H(z),\frac{J_z\eta}{(rz^2+s)^2}\right).}
\end{equation}
Here \(\mathcal B\) and \(J_z\) are \eqref{eq:branch_quartic} and \eqref{eq:J}. Direct substitution proves the equation. More intrinsically,
\begin{equation}\label{eq:vandermonde}
 \chi=p^2(z-v)(z-u)(v-u)
\end{equation}
is invariant under 3-cycles and changes sign under transpositions. The field \(\CC(w,\chi)\) is therefore exactly the fixed field of \(C_3\). This proves degree, kernel, and quotient meaning, rather than just agreement of two equations. The commutative diagram is
\begin{equation}\label{eq:isogeny_diagram}
 \begin{array}{ccc}
 E&\xrightarrow{\ \phi\ (3)\ }&E'=E/C_3\\
 \downarrow 2&&\downarrow 2\\
 \PP^1_z&\xrightarrow{\ H\ (3)\ }&\PP^1_w.
 \end{array}
\end{equation}
Both vertical maps are quotients by inversion with the selected origins. In particular the cubic's four critical values are the images of \(E'[2]\) on the lower target. They are not the four branch points of the left vertical map.

\subsection{Normalization and the modular functions}
On the smooth locus choose \(\sigma^2=s/r\) and set
\begin{equation}\label{eq:normalization}
 z=\sigma Z,\qquad \eta=U\sigma Y_0,\qquad
 w=(p/r)\sigma W,\qquad t=\frac{qr}{ps}.
\end{equation}
The normalized map and curves become
\begin{align}\label{eq:normalized}
 W&=\frac{Z(Z^2+t)}{Z^2+1},\notag\\
 E_t:&\quad Y_0^2=-4Z^4+(t^2-6t-3)Z^2-4t,\notag\\
 E'_t:&\quad \chi_0^2=-4W^4+(t^2+18t-27)W^2-4t^3.
\end{align}
The generic domain is \(t\notin\{0,1,9,\infty\}\). The square-root choice changes the coordinate embedding, not \(t\). With \(M=t^2-6t-3\),
\begin{align}\label{eq:normalized_invariants}
 c_4&=M^2+192t=(t+3)(t^3-15t^2+75t+3),\notag\\
 c_6&=M(576t-M^2),\qquad \Delta=t(t-1)^6(t-9)^2,\\
 j(E_t)&=\frac{(t+3)^3(t^3-15t^2+75t+3)^3}
                   {t(t-1)^6(t-9)^2},\label{eq:jt}\\
 j(E'_t)&=\frac{(t+3)^3(t^3+225t^2-405t+243)^3}
                   {t^3(t-1)^2(t-9)^6}.\label{eq:jprime}
\end{align}
The normalized torsion coordinates are \(P_t=((t+3)^2/12,-(t-1)^2/2)\). To identify them with \(\tau(O)\), choose a critical point \(b\) of the normalized cubic. Then
\begin{equation}\label{eq:torsion_link}
 t=\frac{b^2(b^2+3)}{b^2-1},\qquad a=\frac{2b}{b^2-1}.
\end{equation}
The corresponding ordered fiber is \((a,b,b)\). Its cyclic image is \((b,b,a)\), whose quartic ordinate is \(b(2b^2+3-t)\). Formula~\eqref{eq:weierstrass_map} gives exactly \(P_t\). The exceptional denominators are excluded on the smooth locus or handled in the projective chart. This calculation also explains the orientation convention in \eqref{eq:P}.

Now define
\begin{equation}\label{eq:h}
 h=\frac{t(t-9)^2}{(t-1)^2}.
\end{equation}
Exact substitution gives
\begin{equation}\label{eq:jh}
 \boxed{j(E)=\frac{(h+27)(h+3)^3}{h},\qquad
 j(E')=\frac{(h+27)(h+243)^3}{h^3}.}
\end{equation}
Dualizing the isogeny exchanges these formulas by \(h\mapsto729/h\). They agree with the standard level-three formulas in Elkies' exposition, with this source/quotient orientation recorded \cite[§4, Eq.~(80)]{elkies}.

\subsection{Why the two marked modular curves classify the cubic}\label{sec:moduli_proof}
Here equivalence means independent source and target Möbius changes, with the stated markings preserved. It is not dynamical conjugacy, and it does not preserve the original real plotting coordinate.

\begin{proposition}\label{prop:moduli}
The generic cubic cover without reciprocal marking is classified by the oriented datum \((E,\langle P\rangle)\), an open subset of the coarse modular curve \(X_0(3)\). Retaining reciprocal marking gives \((E,\langle P\rangle,Q)\), equivalently a cyclic subgroup of order six, and an open subset of \(X_0(6)\). The coordinates above are Hauptmoduln: \(h\) on \(X_0(3)\) and \(t\) on \(X_0(6)\).
\end{proposition}
\begin{proof}
The splitting field and its normal order-three subgroup recover \((E,\langle P\rangle)\) from the cover. Conversely any elliptic curve with a cyclic order-three subgroup gives \eqref{eq:isogeny_diagram} by quotienting its 3-isogeny through inversion. This is a degree-three sphere cover with four simple branch values. Choosing a conjugate transposition subgroup changes only the intermediate source isomorphism, not the cover class. The origin can be chosen among the fixed points of that involution; changing it gives an isomorphic datum.

Reciprocity adds the uniquely specified commuting translation \(T_Q\). Conversely, given \((E,\langle P\rangle,Q)\), it descends to nontrivial involutions on both Kummer spheres. The fixed points on the source Kummer sphere are the classes of solutions of \(2R=Q\); there are two such classes. The 3-isogeny is an isomorphism on 4-torsion, and therefore takes this pair bijectively to the corresponding target fixed pair. Normalize the matching pairs to zero and infinity. The cubic becomes an odd map of the form \eqref{eq:odd_cubic}. Independent source and target scalings give \eqref{eq:normalized}. Reversing both matching fixed pairs replaces source and target coordinates by reciprocal coordinates and leaves \(qr/(ps)\) unchanged. No other generic coordinate freedom remains. Thus \(t\) is a complete generic marked coordinate.

There are three choices of nonzero \(Q\) for a generic \((E,\langle P\rangle)\). Forgetting \(Q\) has degree three. The rational function \eqref{eq:h} also has degree three, and the two rational functions in \eqref{eq:jh} generically determine \(h\): over \(\CC(h)\), the gcd of the cleared equations \(j_E(k)=j_E(h)\) and \(j_{E'}(k)=j_{E'}(h)\) is \(k-h\). This exact polynomial check is preserved in the verification suite. Since the ordered pair of \(j\)-invariants is invariant under forgetting \(Q\), \(\CC(h)\) is the fixed moduli field; the equality of degrees shows there is no intermediate missing parameter. Hence \(h\) is the asserted Hauptmodul. The construction also identifies \(\langle P,Q\rangle\) as cyclic of order six, giving \(X_0(6)\).
\end{proof}

A generic elliptic curve has four cyclic subgroups of order three (the four lines of \(E[3]\cong\mathbb F_3^2\)). Accordingly \(h\mapsto j(E)\) has degree four, while \(t\mapsto j(E)\) has degree twelve. Bare \(j\) forgets both the selected 3-isogeny and the reciprocal 2-torsion. Generators \(P\) and \(-P\) give the same cyclic subgroup and are identified by inversion on coarse moduli; an additional rigid labeling of all group elements would be extra data.

The original coefficients satisfy
\begin{equation}\label{eq:ABC_modulus}
 U=(A-C)^2-(B-1)^2,\qquad V=(3A+C)^2-(B+3)^2,
 \qquad t=V/U.
\end{equation}
There is no extra algebraic restriction on generic \(t\). For instance a rational section is
\begin{equation}\label{eq:section}
 (A,B,C)=\left(\frac{t+3}{1-t},\ 1,\ \frac{t-5}{1-t}\right).
\end{equation}
It gives \(p=r=s=8/(1-t)\), \(q=8t/(1-t)\). Thus the three original coefficients contain one intrinsic marked modulus and two generic dimensions of coordinate embedding data. This is a complex-cover classification, not a claim that all corresponding real plots coincide.

\fig{03_moduli_tower}{The cubic sphere map descends from the 3-isogeny. The marked modular coordinate \(t\) retains a nonzero 2-torsion point; \(h\) forgets that choice; \(j(E)\) also forgets the selected cyclic 3-subgroup. Degrees shown are generic coarse degrees.}{fig:moduli}

\subsection{Four critical values and their cross-ratio}
Write the four target branch values as \(c,-c,d,-d\). With a chosen ordering,
\begin{equation}\label{eq:crossratio}
 m'=\left(\frac{d-c}{d+c}\right)^2
\end{equation}
is a Legendre parameter of \(E'\). Reordering gives its six anharmonic transforms. If \(L=t^2+18t-27\), one compatible algebraic expression, up to that reordering, is
\begin{equation}\label{eq:crossratio_t}
 m'=\frac{L+8t^{3/2}}{L-8t^{3/2}},\qquad
 j(E')=256\frac{(1-m'+m'^2)^3}{m'^2(1-m')^2}.
\end{equation}
For comparison, the source double cover \(E\to\PP^1_Z\) has a Legendre parameter
\(m=(M+8\sqrt t)/(M-8\sqrt t)\), with the same ordering caveat. These are algebraic branch choices, not single-valued rational functions of \(t\). The unordered target branch set determines \(j(E')\), but it does not determine the selected incoming 3-isogeny. Four labeled branch values retain additional level-two information.

\section{Exceptional cubic strata and the quartic checkpoint}\label{sec:exceptions}
\subsection{Complete cubic ramification classification}
Restrict first to the degree-three condition \(ps(ps-qr)\ne0\). Define
\begin{equation}\label{eq:K3}
 K_3=3AC+B^2-3B-C^2,\qquad 9ps-qr=-8K_3.
\end{equation}
The critical divisor \eqref{eq:J} and its discriminant give the complete table below. The ramification column lists the index of the unique ramified point over each distinct branch value.
\begin{table}[htbp]\centering\small
\begin{tabularx}{\linewidth}{Yllll}\toprule
Condition & Indices & Monodromy & Deck & Closure genus\\\midrule
\(qr\ne0,\ qr\ne9ps\) & \(2,2,2,2\) & \(S_3\) & trivial & 1\\
Exactly one of \(q,r\) zero & \(3,2,2\) & \(S_3\) & trivial & 0\\
\(q=r=0\) & \(3,3\) & \(C_3\) & \(C_3\) & 0\\
\(qr=9ps\) & \(3,3\) & \(C_3\) & \(C_3\) & 0\\\bottomrule
\end{tabularx}
\caption{All degree-three cases. Closure degrees are six for the \(S_3\) rows and three for the \(C_3\) rows.}\label{tab:cubic_strata}
\end{table}

\begin{proof}
A degree-three sphere map has total ramification four and local degree at most three. Its critical divisor can therefore consist only of four simple points, one double point and two simple points, or two double points. Distinct critical points cannot share an image: even two simple ramification points would use degree four in one fiber. Formula~\eqref{eq:J} gives exactly the rows displayed, including infinity through its binary form. In the first two rows, transitivity and the presence of a transposition force \(S_3\); in the last two, the branch cycles generate \(C_3\). Galois Riemann--Hurwitz gives respectively
\begin{equation}\label{eq:RH_strata}
 \begin{array}{ll}
 6[-2+4(1-1/2)]=0,&g=1,\\
 6[-2+(1-1/3)+2(1-1/2)]=-2,&g=0,\\
 3[-2+2(1-1/3)]=-2,&g=0.
 \end{array}
\end{equation}
No critical-value collision locus has been omitted within the degree-three domain.
\end{proof}

The residual quadratic splits over \(\CC(z)\) exactly on the two cyclic rows:
\begin{equation}\label{eq:cyclic_loci}
 q=r=0\quad\text{or}\quad qr=9ps.
\end{equation}
Indeed a rational root would give a nontrivial deck transformation and hence a cyclic cover; conversely a cyclic cover has two rational partners. In original coefficients these loci are
\begin{equation}\label{eq:cyclic_ABC}
 B=-3,\quad C=-3A,\qquad\text{or}\qquad K_3=0,
\end{equation}
always with \eqref{eq:degree3}. They are disjoint there. On the first, \(w=(p/s)z^3\), with deck maps \(z\mapsto\omega z\), \(\omega^3=1\). On the second, choose \(a^2=3s/r\), put \(\kappa=3pa/r\), and observe
\begin{equation}\label{eq:cyclic_conjugacy}
 \frac{w-\kappa}{w+\kappa}=\left(\frac{z-a}{z+a}\right)^3.
\end{equation}
Conjugating multiplication by \(\omega\) gives the Möbius partners explicitly. These formulas are over \(\CC\), without a claim that all such partners are real for real coefficients.

Adding reciprocal equivariance to the cyclic deck group produces \(C_6\) on \(q=r=0\), since \(-z\) commutes with multiplication by \(\omega\). On \(qr=9ps\) it produces the dihedral group of order six, since \(({-z}-a)/({-z}+a)\) is the reciprocal of \((z-a)/(z+a)\). These extended groups do not fix the target pointwise. Generically the only source transformations fixing or inverting the value are the identity and reciprocity: composing any second value-inverting map with reciprocity would give a deck transformation.

\subsection{Cancellation and singular limits are different operations}
The lower-degree strata follow directly from the homogeneous pair after \eqref{eq:pqrs}. Constants are \(A=1,B=C\) for \(+1\), and \(A=-1,C=-B\) for \(-1\). A nonconstant point has degree two exactly when one of \(p,s\) is zero and \(ps-qr\ne0\). Every other nonconstant point on the resultant-zero locus has degree one. A factor at zero comes from \(s=0\), a factor at infinity from \(p=0\), and the reciprocal quadratic factor from \(ps=qr\) when \(p,s\ne0\). Their intersections give the stated degree-one or constant cases. Covers of degree two and one have respectively \(C_2\) and trivial monodromy, with genus-zero closures.

Degenerating the quartic model \(E\) is not the same operation as taking the connected Galois closure of the specialized rational map. On \(qr=9ps\), for example,
\begin{equation}\label{eq:degenerate_square}
 \mathcal D(z)=-\frac{4p^2s}{r}(rz^2-3s)^2,
\end{equation}
so the compactified correspondence splits into two rational components meeting at two ordinary nodes. Meanwhile
\begin{equation}\label{eq:branch_square}
 \mathcal B(w)=-4r^3s\left(w^2-\frac{27p^2s}{r^3}\right)^2,
\end{equation}
and the four simple critical values coalesce in pairs. The specialized cubic itself is a smooth cyclic cover with connected closure \(\PP^1\); the reducible limit of the old degree-six closure is a different object.

If \(q=0,r\ne0\), then
\(\mathcal D=-p^2s z^2(4rz^2+3s)\).
It has an ordinary node at zero and irreducible rational normalization, still with \(S_3\) monodromy. If \(r=0,q\ne0\), the analogous double root is at infinity. If \(q=r=0\), the binary quartic is \(-3p^2s^2Z^2T^2\), giving two rational components with nodes at zero and infinity. The genus-zero correspondence in the one-zero cases remains irreducible over the source field and does not create a rational partner.

\subsection{The four modular boundary points}
In the normalized family the excluded values are \(t=0,1,9,\infty\). The Weierstrass discriminant has orders \(1,6,2,3\), respectively; the last follows by the minimal scaling \(X=t^2\bar X\), \(Y=t^3\bar Y\) near infinity. In each case the correspondingly scaled \(c_4\) is nonzero. These are multiplicative, nodal degenerations of the elliptic model, not cuspidal cubic degenerations. For the associated semistable elliptic surface, the minimal resolved fibers have types \(I_1,I_6,I_2,I_3\). One should distinguish that resolved surface model from an uncontracted incidence model such as \eqref{eq:degenerate_square}.

The modular maps make the boundary orders transparent:
\begin{equation}\label{eq:h_ramification}
 \frac{dh}{dt}=\frac{(t-9)(t+3)^2}{(t-1)^3},\qquad
 h+27=\frac{(t+3)^3}{(t-1)^2}.
\end{equation}
The two cusps of \(X_0(3)\) are \(h=0,\infty\), with source-\(j\) pole orders one and three. Their preimages give the four \(X_0(6)\) cusp widths listed above. At \(h=-27\), the selected level-three datum has an order-three elliptic stabilizer; its threefold preimage is \(t=-3\). The zero at \(h=-3\) is instead triple in \(j(E)\). Also
\begin{equation}\label{eq:j1728}
 j(E)-1728=\frac{(h^2+18h-27)^2}{h}.
\end{equation}
There is no order-two elliptic point on \(X_0(3)\); the two roots in \eqref{eq:j1728} have local degree two over the bare elliptic invariant. Retaining \(Q\) eliminates these extra elliptic stabilizers beyond the generic inversion on the coarse \(X_0(6)\) datum.

At intersections of coefficient strata, the normalized chart can itself be singular. For example \(q=r=0\) has \(t=0\), but a path along which both vanish simply makes \(t\) vanish to order two, while \(\sigma^2=s/r\) diverges. The two-component incidence limit there does not contradict the simple nodal fiber of the normalized \(t\)-model at zero. The paper makes no exhaustive claim about all boundary paths or their chosen total-space resolutions.

\fig{04_boundaries}{The degree-three parameter strata separate smooth elliptic closure, rational \(S_3\) normalization, cyclic splitting, and degree loss. The normalized coordinate reaches the four cusps with the displayed source-\(j\) pole orders. Coefficient intersections can require different local charts.}{fig:boundaries}

\subsection{Quartic checkpoint}\label{sec:quartic}
For
\begin{equation}\label{eq:quartic}
 F_4(y)=\frac{Ay^4-By^3-Cy^2-Dy+1}{y^4-Dy^3-Cy^2-By+A},
\end{equation}
the exact fixed-value numerators are
\begin{align}\label{eq:quartic_fibers}
 P-Q&=(y^2-1)[(A-1)(y^2+1)+(D-B)y],\notag\\
 P+Q&=(A+1)(y^4+1)-(B+D)(y^3+y)-2Cy^2.
\end{align}
Set
\begin{equation}\label{eq:H4}
 \begin{split}
 H_4={}&A^3+A^2C-A^2+ABD-2AC-AD^2\\
       &-A-B^2+BD+C+1.
 \end{split}
\end{equation}
Then
\begin{equation}\label{eq:quartic_resultant}
 \Res(P,Q)=(A-B-C-D+1)(A+B-C+D+1)H_4^2.
\end{equation}
The Wronskian is
\begin{equation}\label{eq:W4}
 W_4=a_4(y^6+1)+b_4(y^5+y)+c_4^{\rm crit}(y^4+y^2)+d_4y^3,
\end{equation}
with
\begin{equation}\label{eq:W4_coeff}
 \begin{aligned}
 a_4&=B-AD,& b_4&=2C(1-A),\\
 c_4^{\rm crit}&=-3AB+BC-CD+3D,&d_4&=4A^2+2B^2-2D^2-4.
 \end{aligned}
\end{equation}
The superscript distinguishes this critical-polynomial coefficient from an elliptic invariant. With \(\zeta=y+1/y\), its finite nonzero roots are obtained from
\begin{equation}\label{eq:critical_cubic4}
 a_4\zeta^3+b_4\zeta^2+(c_4^{\rm crit}-3a_4)\zeta+d_4-2b_4=0,
 \qquad y^2-\zeta y+1=0.
\end{equation}
On the generic locus there are six simple, distinct branch values. The closure has group \(S_4\), degree 24, and
\begin{equation}\label{eq:quartic_genus}
 2g-2=24[-2+6/2]=24,\qquad g=13.
\end{equation}
There is no generic deck transformation; the residual partner relation is an irreducible cubic. The return of a \(+1\) lock at \(-1\) is an even-degree effect, not a return of quadratic Galois geometry.

Three previously established exceptional families suffice to demonstrate the limits of the generic theorem. In each example assume the displayed map remains degree four.
\begin{enumerate}
\item \(B=D=0\): \(F_4(y)=F_{2;A,C}(y^2)\), a decomposable cover with deck symmetry \(y\mapsto-y\).
\item \(B=C=D=0\), \(A^2\ne1\): \(F_4\) is a target Möbius transform of \(y^4\), a cyclic \(C_4\) Galois cover with genus-zero source.
\item \(A=-1\), \(B=D=0\), \(C\ne0\): the deck maps are \(y,-y,i/y,-i/y\). They form \(V_4\), and the cover is Galois of genus zero.
\end{enumerate}
These are selected loci, not an exhaustive quartic census.

\subsection{The squaring lift and real plots}\label{sec:x_lift}
For a reduced degree-\(d\) map, \(f_n(x)=F_n(x^2)\) has degree \(2d\) and obeys
\begin{equation}\label{eq:x_lift}
 f_n(-x)=f_n(x),\qquad f_n(1/x)=1/f_n(x),\qquad
 e_{f_n}(x)=e_{x^2}(x)e_{F_n}(x^2).
\end{equation}
The first is a deck symmetry; the second is equivariance. For generic coefficients, each old simple branch value has two simple ramification points in the lift, with partition \(2^2 1^{2n-4}\). Zero and infinity of the \(x\)-sphere add branch values \(1/a_0\) and \(a_0\), each with partition \(2,1^{2n-2}\). These values are generically distinct from one another and from the old branch values. Thus the cubic lift has degree six and six branch values; the quartic lift has degree eight and eight branch values. Total ramification is \(2(2n-2)+2=4n-2\), as required. The genus formula in Theorem~\ref{thm:generic} is not a formula for these lifted covers or their further target quotients.

For real plotting, only \(y=x^2\ge0\) is sampled. Positive zeros or poles give paired real locations \(x=\pm\sqrt y\); negative real locations in \(y\) give imaginary \(x\), and nonreal roots have complex square roots. Generic locks at \(y=1\) give \(x=\pm1\). The parity-controlled behavior at \(y=-1\) lies at \(x=\pm i\), outside the real plot. General complex Möbius equivalences do not preserve this half-line or the original embedding.

Even real coefficients need not make the correspondence an elliptic curve with a real origin. At \(t=2\), realized for instance by \((A,B,C)=(-5,1,3)\),
\begin{equation}\label{eq:real_torsor}
 E_2:\ Y_0^2=-4Z^4-11Z^2-8
\end{equation}
has no real points, including at infinity. Its Jacobian has a real identity, but this genus-one curve is a real torsor without a real point. The complex Weierstrass transformation selected a complex origin. Real twists, branch ordering, and the \(y\ge0\) restriction are extra data not captured by complex \(j\) or \(t\).

\section{Final bounded Tri-Octagon bridge test}\label{sec:bridge}
This section is a comparison with pre-existing project definitions, not a physical interpretation or a proposed change to the shell. The six objects being compared are \emph{candidate measurement patches}. The source explicitly states that the welded shell has two connected nine-edge rims, not six boundary components \cite[§§1--2]{patch}. Their group action was derived in the external cubic report \cite[§12]{cubic}. We preserve that result and ask only whether the already defined continuous quantities supply a modular identification.

\subsection{The exact finite correspondence}
Use fresh geometric notation to avoid confusing a patch parameter with the modular coordinate \(t\). Put
\begin{equation}\label{eq:patch_constants}
 s_g=\sqrt2-1,\quad d=1-\sqrt2/2,\quad
 O=(0,\sqrt3/6,0),\quad h_0=s_g/2,\quad r_0=1/\sqrt3,
\end{equation}
and \(u_j=R_z(2\pi j/3)(0,1,0)\), \(v_j=R_z(2\pi j/3)(-1,0,0)\). The existing finite triangles are
\begin{equation}\label{eq:patches}
 \begin{split}
 \mathcal P_j^\epsilon:\quad
 X&=O+r(\ell)u_j+\omega v_j+\epsilon h(\ell)e_z,\\
 r(\ell)&=r_0-\sqrt3d\ell/2,\qquad h(\ell)=h_0+d\ell,\\
 0&\le\ell\le1,\quad |\omega|\le d\ell/2,
 \quad j\in\ZZ/3,\quad\epsilon\in\{1,-1\}.
 \end{split}
\end{equation}
On coordinates centered at \(O\), define the actual rigid maps
\begin{equation}\label{eq:patch_group}
 R=R_z(2\pi/3),\qquad S=\operatorname{diag}(-1,1,-1).
\end{equation}
Then \(R^3=S^2=1\), \(SRS=R^{-1}\). They act on whole patches by
\begin{equation}\label{eq:patch_action}
 R:\mathcal P_j^\epsilon\mapsto\mathcal P_{j+1}^\epsilon,
 \qquad S:\mathcal P_j^\epsilon\mapsto\mathcal P_{-j}^{-\epsilon}.
\end{equation}
Indeed \(Su_j=u_{-j}\), \(Sv_j=-v_{-j}\), \(Se_z=-e_z\); the change \(\omega\mapsto-\omega\) preserves the finite triangle inequalities. No nonidentity rotation fixes a patch label, and every involution swaps upper and lower. The resulting \(S_3\) action is free and transitive.

The explicit bijection
\begin{equation}\label{eq:finite_bridge}
 \boxed{\Phi:\mathcal P_j^\epsilon\longmapsto(j,j+\epsilon)}
\end{equation}
identifies the patches with the six ordered pairs of distinct elements of \(\ZZ/3\). Under \(R\), both entries increase by one; under \(S\), both are negated. Forgetting the second entry corresponds to forgetting upper/lower and retaining the vertical pair \(j\). Its three stabilizers are the three transposition subgroups. This is an exact equivalence of finite group actions and incidence projections, stronger than a match of the numbers three and six.

It matches the left regular monodromy action on root labels. The deck action \eqref{eq:generators} permutes positions and gives the commuting right regular action on complete root orderings. Identifying these actions without that distinction would obscure which quotient is being used. The bijection also requires a seed label and orientation; it is not a canonical assignment of spatial patches to analytic branches.

The pure upper/lower reflection at fixed \(j\) commutes with \(R\), producing \(C_3\times C_2\), not \(S_3\). The noncommuting involution in \eqref{eq:patch_group} is specifically a half-turn that reverses pair index as well as height. This geometric step is essential to the established finite result.

\subsection{What would constitute a continuous bridge?}
The candidate targets are not interchangeable. The coordinate \(t\) determines a reciprocal-marked cubic cover; \(h\) determines its oriented cyclic 3-isogeny; and an ordered four-critical-value cross-ratio determines a labeling of the branch divisor of the quotient elliptic curve \(E'\). A structure-preserving identification must derive the relevant data from the shell definitions. It must not choose a convenient scalar and declare it equal to a modular coordinate.

In particular, a cross-ratio requires four independently selected distinct points on a common projective line, with a stated ordering or its permutation class. For \(h\) one needs the selected 3-isogeny, not merely an elliptic \(j\). For \(t\) one additionally needs the selected nonzero 2-torsion. The finite regular \(S_3\)-set supplies none of these continuous choices by itself.

There is a precise non-identifiability argument. For every \(t\) in the smooth domain of \eqref{eq:normalized}, the regular fiber has exactly the same finite \(S_3\)-set and the same three-sheet incidence projection. Therefore \eqref{eq:finite_bridge} can be made at every such \(t\). Two generic distinct \(t\)'s define inequivalent reciprocal-marked covers by Proposition~\ref{prop:moduli}, yet their finite incidence data are isomorphic. Thus the finite bridge cannot determine \(t\), \(h\), or the four-point branch configuration.

\subsection{Audit of independently defined continuous quantities}
The bounded test consulted the exact definitions below, preserved as numbered snapshots with hashes. The full source map appears in Appendix~\ref{app:bridge_sources} and the companion ledger. No kernel module was imported or executed for this comparison.

\paragraph{Fixed welded-shell shape.}
The patch source defines normals
\(n_j^\epsilon=(2u_j+\epsilon\sqrt3e_z)/\sqrt7\),
spatial side lengths \((s_g,s_g,d)\), tip angle \(\arccos(3/4)\), and projected equilateral triangles of side \(d\). It also defines tip, centroid, and base-midpoint radius/height ratios. For example the centroid has
\begin{equation}\label{eq:fixed_ratio}
 r_c/h_c=2\sqrt3-\sqrt6.
\end{equation}
These dimensionless values are constants for the accepted shell up to similarity. An invariant of this fixed similarity class cannot furnish a nonconstant modular family. A chosen constant could be assigned to one elliptic curve, but the definitions select neither that curve nor an isogeny. The parameter \(\ell\) in \eqref{eq:patches} runs over points inside a fixed triangle; it is not a deformation modulus of the shell.

\paragraph{Separate reference scaffold.}
The pre-existing reference scaffold \cite{scaffold} has independent positive lengths \(s\) and \(g_{\rm gap}\), alternating side lengths \((s,g_{\rm gap})\), and
\begin{equation}\label{eq:scaffold}
 \begin{split}
 g_{\rm gap}&=\sqrt3p-s/2,\qquad
 R_H^2=(s^2+sg_{\rm gap}+g_{\rm gap}^2)/3,\\
 \mathcal A_H&=\sqrt3(s^2+4sg_{\rm gap}+g_{\rm gap}^2)/4.
 \end{split}
\end{equation}
Modulo scale, \(g_{\rm gap}/s\) is a genuine shape parameter, with regularity at one. The selected edge/connector roles are retained by its class-preserving \(D_3\) symmetry, so that ratio is invariant under that action. Paper C's member has \(g_{\rm gap}/s=1/\sqrt2\); varying the scaffold is a separate family, not a hidden deformation already present in the fixed welded shell.

Neither the defining formulas nor the documented derivation supplies a projective four-point branch divisor, an elliptic group law, or a selected isogeny. Both a shape ratio and \(t\) being invariant under a finite symmetry is compatible with many arbitrary functions between them; it does not single out one. Defining \(t=g_{\rm gap}/s\), or applying a rational reparameterization to force the cusp values, would add precisely the dictionary being tested. We do not make that assignment. The natural six scaffold vertices are also not already four critical values; selecting four would require a new rule and branch-cover interpretation.

\paragraph{Orientation, readout, and conditional registration.}
The source harmonic is \(A\cos3(\theta-\ell_0)\). Its six extrema have three vertical pairs, while the fixed sampled clock is a different finite set. The source's finite-patch registration uses an additionally assumed coordinate dictionary and a free relative angle \(\delta\), with a linked scale
\begin{equation}\label{eq:registration}
 a(\delta)=\frac{K}{\cos\delta+\sqrt3/2},\qquad
 K=\sqrt3/12+\sqrt6/4.
\end{equation}
That parameter chooses an attachment of an observer construction to the shell; it is not an intrinsic shell invariant. The report leaves a unique gate anchor, axis, and channel assignment undetermined. Rotations or a choice of readout origin do not produce four projectively defined branch points. Treating one of these angles as a modular coordinate would add an attachment convention and an analytic structure absent from the accepted definitions.

The existing readout code also defines \(J_{\rm read}=\operatorname{Im}(\Omega_1\overline{\Omega_2}\Omega_3)\), its bounded normalization \(J_{\rm read}/(1+|J_{\rm read}|)\), and a memory update \cite{readout}. The function name \texttt{cubic\_j} does not make it an elliptic \(j\)-invariant. It is a state-dependent real cubic expression and is not even invariant under a cyclic permutation of its arguments: \((1,i,1)\) gives \(-1\), while \((i,1,1)\) gives \(+1\). No documented transformation law turns this readout into one of the cover invariants.

\paragraph{Existing boundary response and transfer.}
The boundary-response module explicitly labels its choice a toy response \cite{boundary}. Independently of the Three-Way hierarchy, it defines
\begin{equation}\label{eq:lens}
 \theta=\arccos\frac{d_\ell}{2r_\ell},\qquad
 I(\theta)=\frac{2\theta-\sin2\theta}{\pi},\qquad
 \mathcal P(\theta,\xi)=\sqrt{I(\theta)}(\xi\otimes f_0),
\end{equation}
with \(0\le\theta\le\pi/2\) and \(f_0=(1,1,1)/\sqrt3\). The area fraction ranges from zero to one and is strictly increasing on the interior, since \(I'(\theta)=4\sin^2\theta/\pi>0\). It is a lens geometry scalar, and \(\xi\) is an independent incident two-vector. Neither specifies a cubic cover or branch labeling. The interval is not itself an obstruction to parameterizing a real modular arc; the missing point is an independently defined structure-preserving rule.

The existing SRG module fixes numerical parameters and a transfer \(U=B\otimes(RZC)\), with a specified two-by-two helicity matrix and chosen eigenmode gauge \cite{srg}. Its transfer count is discrete and incident state is externally supplied. Its fixed spectra and projectors supply no documented varying four-point branch divisor or selected elliptic isogeny. Selecting four among eigenvalues, or applying an arbitrary transform to an input amplitude, would be a new construction. The operating-region profile bounds this initialization and an existing recurrence; those bounds do not introduce modular structure \cite{operating}. These definitions were read for the bridge audit only, without changing or retesting the production implementation.

\subsection{Bridge classification and endpoint}\label{sec:bridge_result}
\begin{center}
\fbox{\parbox{.9\linewidth}{\centering\textbf{FINITE \(S_3\) BRIDGE ONLY}\par
The patch action and its three-pair quotient match exactly. No continuous modular identification is independently justified by the inspected definitions.}}
\end{center}
The positive result is \eqref{eq:finite_bridge}. The fixed shell has no varying similarity modulus. The separate scaffold and response systems do have continuous inputs, but their definitions do not give the missing elliptic or branch-cover data. Consequently this bounded test establishes no structure-preserving map to \(t\), \(h\), or the four-critical-value cross-ratio.

This is a source-limited negative result, not a theorem that no future mathematical construction could connect these subjects. The continuous part is underdetermined by the inspected definitions; the overall classification is ``finite bridge only'' because the finite equivalence is already proved. No scalar was fitted, no new model was introduced, and no physics claim follows. The comparison stops at this endpoint.

\fig{05_finite_bridge}{The preserved finite correspondence intertwines actual rigid patch symmetries with ordered-pair label permutations. Every smooth cubic parameter has this same finite incidence system. The missing continuous data are a branch divisor and its covering or isogeny marking; the figure asserts no map from shell dimensions to a modular coordinate.}{fig:bridge}

\section{Closeout}\label{sec:closeout}
Coefficient reversal preserves equivariance, reciprocal divisors, parity-sensitive fibers, and an odd Cayley normal form throughout the hierarchy. It does not preserve the quadratic deck partner. The generic other-sheet relation becomes irreducible, symmetric monodromy controls the closure, and Riemann--Hurwitz gives the exact genus formula. The cubic is the first elliptic level of this particular extension, with an explicit translation/inversion action and a 3-isogeny. The quartic already has generic closure genus thirteen.

The intrinsic cubic modulus is one-dimensional after coordinate embeddings are forgotten, with the selected marking distinguishing \(X_0(6)\) from \(X_0(3)\) and from bare elliptic \(j\). This geometry is established independently of the Tri-Octagon comparison. The latter ends with an exact finite incidence match and no justified continuous extension.
\begin{equation*}
 \boxed{\text{THREE-WAY RECIPROCAL HIERARCHY = PARKED}}
\end{equation*}
The paper undertakes no further degree-five study, quartic census, physical interpretation, kernel integration, or new shell coupling.

\clearpage\appendix\section{Verification, evidence classes, and reproducibility}\label{app:evidence}
\subsection{Evidence hierarchy}
The general monodromy and genus result is a theorem with an existence proof in Section~\ref{sec:generic}. Its topology is not claimed to be machine-formalized. The displayed identities, discriminants, birational maps, torsion formulas, and isogeny formulas have exact symbolic checks. Numerical evaluations are additional consistency tests, not substitutes for those proofs. The final bridge classification is a bounded audit of specified definitions together with an exact finite-action equivalence; it is not a universal impossibility theorem.

The companion \texttt{provenance/claim\_ledger.csv} maps every theorem and major derived claim to a source-report section, manuscript label, verification item or prose proof, and limitation. Standard facts, explicit family derivations, and project-specific source comparisons are separately marked. No mathematical priority is claimed. The literature check was bounded to primary sources needed for normal-form and modular conventions; it was not an exhaustive novelty search.

\subsection{Exact suites and numerical witnesses}
The original external verification scripts were copied byte-for-byte into this paper's independent verification directory and rerun there. Their original reports, results, and scripts remain unchanged. With SymPy 1.14.0 and mpmath 1.3.0:
\begin{itemize}
\item The hierarchy suite passed 51 exact checks and four numerical checks, covering degrees two through four, cancellation factors, critical equations, cubic splitting, and exceptional examples.
\item The cubic suite passed 73 exact checks, including quartic/Weierstrass identities, 3-torsion, the selected 2-torsion, root-incidence actions, the 3-isogeny, modular functions and generic recovery of \(h\), degeneration formulas, and the finite patch action.
\item Ten additional exact checks verify the independently sourced geometric ratios, lens-area derivative and endpoints, and the cyclic readout counterexample used in the final audit. These checks introduce no fitted dictionary.
\end{itemize}
This is 134 exact checks in total. The two older suites retain their original test names and output structures. The extra checks are isolated in \texttt{verification/verify\_bridge.py}.

Two fixed coefficient witnesses certify the nonempty simple-branch loci in the low degrees. For \((A,B,C)=(2,3,5)\),
\begin{align}\label{eq:witness3}
 W_3&=-7y^4-2y^3-7y^2-2y-7,\notag\\
 \Disc(W_3)&=17000000,\notag\\
 \Disc_y(P-vQ)&=5(353v^4-1222v^3+1763v^2-1222v+353).
\end{align}
The last polynomial has discriminant \(1228250000000000\ne0\). For \((A,B,C,D)=(2,3,5,7)\),
\begin{align}\label{eq:witness4}
 W_4&=-11y^6-10y^5-17y^4-68y^3-17y^2-10y-11,\notag\\
 \Disc(W_4)&=376072090731675648,\notag\\
 \Disc_y(P-vQ)&=-96(7151v^6-35048v^5+78095v^4\notag\\
 &\hspace{24mm}-100468v^3+78095v^2-35048v+7151).
\end{align}
Its discriminant is the nonzero integer recorded in the hierarchy JSON. Both witnesses also have nonzero numerator--denominator resultants. These exact nonvanishings prove simple critical divisors and distinct branch values; a numerical root plot is not needed for that inference.

Numerical evaluation used 100 decimal digits, with critical roots approximated at 85 digits. The maximum reciprocity residuals were \(2.94\times10^{-101}\) and \(1.84\times10^{-101}\) for the cubic and quartic, respectively. Their maximum normalized critical-polynomial residuals were \(6.02\times10^{-86}\) and \(5.48\times10^{-86}\). For three fixed complex probes at \(t=17\), the cubic suite obtained maximum absolute residuals \(2.52\times10^{-95}\) for the Weierstrass map and \(4.09\times10^{-96}\) for the isogeny. These are high-precision floating evaluations, not interval-certified enclosures. There was no parameter scan.

\subsection{Exact source map for the final bridge}\label{app:bridge_sources}
All line numbers below refer to the preserved snapshots of the pre-existing files, not to new definitions in this paper. The source manifest records absolute original paths and SHA-256 hashes. A numbered UTF-8 copy beside each snapshot makes the locators independent of an editor's line display.
\begin{longtable}{L{.18\linewidth}L{.24\linewidth}L{.49\linewidth}}
\toprule ID & Locator & Definition used and outcome\\\midrule\endfirsthead
\toprule ID & Locator & Definition used and outcome\\\midrule\endhead
PATCH & Report §§2--3.1; lines 51--150 & Finite patch coordinates, normals, metric ratios, and rigid symmetries. These support the exact regular \(S_3\)-set, not a branch divisor.\\
PATCH & §§3.3--3.4; lines 161--233 & Conditional registration parameter and frozen clock distinction. These do not select a shell modulus.\\
PATCH & §7; lines 386--393 & Attachment, gate axis, and channel meaning remain source-limited.\\
SCAFFOLD & \texttt{ReferenceScaffold}; lines 172--266 & Positive side/gap family, radius and area, role labels, regularity. A shape modulus exists, but no elliptic or isogeny datum is defined.\\
SCAFFOLD & \texttt{paper\_c\_member}; lines 339--342 & The fixed welded-member placement selects gap ratio \(1/\sqrt2\).\\
SCAFFOLD\_\allowbreak PROOF & §6 & Independent geometric derivation of the alternating-side family; distinguishes regularity from the fixed welded member.\\
BOUNDARY & Lines 1--7, 25--37, 55--105 & Explicit toy lens-area response, domain, gain, and six-vector preparation. No cubic branched cover is specified.\\
SRG & Lines 21--29, 50--94, 122 onward & Fixed parameters, tensor transfer, eigenmode gauge, and discrete transfer count. No continuous modular correspondence is defined.\\
OPERATING & Lines 1--16, 24--31, 46 onward & Bounded initialization/recurrence profile; not a geometric branch parameter.\\
Z\_READOUT & Lines 91--151, 226--260 & State norm, harmonic, cubic state readout, and memory update. The function \texttt{cubic\_j} is not an elliptic invariant.\\\bottomrule
\end{longtable}
The original paths of PATCH and the code snapshots are relative to the recorded repository root; the separate scaffold proof is in the external research tree. The manuscript bibliography supplies their identities. These snapshots serve as evidence, not replacement production modules.

\subsection{Bundle and build contract}
The independent paper directory is \path{papers/three_way_reciprocal_hierarchy}. It contains multi-file LaTeX source, bibliography, vector and raster figures, the verification suites and outputs, source snapshots, claim/source ledgers, and build/audit records. The final PDF is in \path{output/pdf/}. A ZIP bundle preserves the same tree without recursively including itself.

The supplied build script uses the existing Tectonic toolchain and a dedicated external build/cache workspace. It does not install a TeX system or modify Paper I. Verification scripts use the existing mathematical Python environment. Figure generation is deterministic and uses exact formulas or explanatory diagrams; figures are not evidence for a theorem beyond the equations cited in their captions. The PDF is rendered for visual review, and internal destinations, citations, source labels, and figure inclusion are audited.

The baseline records all tracked repository file identities, all Paper I, kernel, and UI files present in the protected directories, and both external report directories. Final validation compares their hashes and repository HEAD; new files are confined to the Paper II directory and external build/QA workspace. Existing user changes are preserved. The detailed build report and artifact manifest accompany the delivery.

\clearpage\phantomsection\addcontentsline{toc}{section}{References}
\begin{thebibliography}{99}
\raggedright
\bibitem{paper1} \emph{Three-Way Reciprocal Geometry: Reconstruction of a Historical Rational Model from Belyi Quotients to an Edwards Elliptic Galois Closure}. Reconstruction draft prepared for Hilmir F. Halldórsson with Codex assistance, v0.1, 5 October 2026. Project manuscript, sibling directory \path{papers/three_way_reconstruction}. Preserved unchanged.
\bibitem{hierarchy} \emph{Three-Way reciprocal hierarchy, n = 2--4: first extension v0.1}. External exact-analysis report, 5 October 2026. Snapshot HIERARCHY; original directory \path{research/three_way_reciprocal_hierarchy_20261005} in the workspace research tree. Includes \path{verify_hierarchy.py} and exact results.
\bibitem{cubic} \emph{Three-Way cubic elliptic correspondence v0.1}. External exact-analysis report, 5 October 2026. Snapshot CUBIC; original directory \path{research/three_way_cubic_elliptic_20261005} in the workspace research tree. Includes \path{verify_cubic_elliptic.py} and exact results.
\bibitem{msw} N. Miasnikov, B. Stout, and P. Williams, \emph{Automorphism loci for the moduli space of rational maps}, arXiv:1408.5655v2 (2014), §2.1, Lemma 2. \url{https://arxiv.org/pdf/1408.5655}. Commuting automorphisms are distinct from deck transformations.
\bibitem{pakovich} F. Pakovich, \emph{On rational functions whose normalization has genus zero or one}, arXiv:1609.03482v3 (2017). \url{https://arxiv.org/pdf/1609.03482}.
\bibitem{elkies} N. D. Elkies, \emph{Elliptic and modular curves over finite fields and related computational issues}, primary exposition (1997), §4, Eq.~(80). \url{https://people.math.harvard.edu/~elkies/modular.pdf}.
\bibitem{patch} \emph{Six candidate apertures, the Z harmonic and a rotated toroidal motif}. Codex research report, RESEARCH\_GATE\_TORUS\_v0.1, 24 September 2026. Repository path \path{research/GATE_TORUS_INVESTIGATION_v0.1/REPORT.md}; snapshot PATCH. Only the existing geometric and registration definitions are used here.
\bibitem{scaffold} Tri-Octagon project, \path{kernel_physics/reference_scaffold.py}, preserved snapshot SCAFFOLD, and \emph{Tri-Octagon hexagon core derivation v0.1}, snapshot SCAFFOLD\_PROOF. Exact paths, hashes, and line locators appear in the source manifest and Appendix~\ref{app:bridge_sources}.
\bibitem{boundary} Tri-Octagon project, \path{kernel_physics/boundary_response.py}, preserved snapshot BOUNDARY. Explicit \path{lens_area_norm_v1} toy response; read only.
\bibitem{srg} Tri-Octagon project, \path{kernel_physics/srg.py}, preserved snapshot SRG. Fixed November transfer and initialization handoff; read only.
\bibitem{operating} Tri-Octagon project, \path{kernel_physics/operating_region.py}, preserved snapshot OPERATING. Bounded initialization and existing recurrence adapters; read only.
\bibitem{readout} Tri-Octagon project, \path{kernel_physics/z_manifold.py}, preserved snapshot Z\_READOUT. Passive harmonic and cubic state readouts; read only.
\end{thebibliography}

\end{document}
